\clearpage
\VerdierPhysicalPageAnchor{0199}%
\begingroup\small\setlength{\parskip}{0pt}\setlength{\emergencystretch}{1.5em}
\setlength{\abovedisplayskip}{4pt plus 2pt minus 2pt}%
\setlength{\belowdisplayskip}{4pt plus 2pt minus 2pt}%

\VerdierSemanticAnchor{ch3d-2-3-2-exact-sequences}
et des suites exactes :
\[
0\longrightarrow H^{-q}(X_{q-1})\longrightarrow
\frac{Z^{-q}}{\operatorname{Im}(Y^{-q-1}\longrightarrow Z^{-q})}
\longrightarrow
\frac{Z^{-q}}{\operatorname{Ker}(Z^{-q}\longrightarrow L^{-q+1})}
\longrightarrow0,
\]
\[
0\longrightarrow
\frac{Z^{-q}}{\operatorname{Ker}(Z^{-q}\longrightarrow L^{-q+1})}
\longrightarrow L^{-q+1}\longrightarrow
\frac{L^{-q+1}}{\operatorname{Im}(Z^{-q}\longrightarrow L^{-q+1})}
\longrightarrow0,
\]
\[
0\longrightarrow H^{-q+1}(X_{q-1})\longrightarrow
\frac{L^{-q+1}}{\operatorname{Im}(Z^{-q}\longrightarrow L^{-q+1})}
\longrightarrow
\frac{L^{-q+1}}{\operatorname{Ker}(L^{-q+1}\longrightarrow L^{-q+2})}
\longrightarrow0.
\]

\VerdierSemanticAnchor{ch3d-2-3-2-mq-object}
Le complexe $X_{q-1}$ étant quasi-isomorphe au complexe $Y$, les objets
$H^i(X_{q-1})$ sont des éléments de $M$. Il résulte alors des suites exactes
écrites ci-dessus, de l'hypothèse de récurrence 2), et des propriétés
d'exactitude de l'ensemble $M^\wedge$ (2.2.4), que l'objet
\[
Z^{-q}/\operatorname{Im}(Y^{-q-1}\longrightarrow Z^{-q})
\]
est un élément de $M^\wedge$. Par suite, il existe un morphisme :

\VerdierSemanticAnchor{ch3d-2-3-2-v}
\[
v\colon L^{-q}\longrightarrow Z^{-q}
\]
tel que $L^{-q}$ soit un élément de $M$ et que le morphisme composé :

\VerdierSemanticAnchor{ch3d-2-3-2-composite}
\[
L^{-q}\xrightarrow{v}Z^{-q}\longrightarrow
Z^{-q}/\operatorname{Im}(Y^{-q-1}\longrightarrow Z^{-q})
\]
soit un épimorphisme. Posons alors :

\VerdierSemanticAnchor{ch3d-2-3-2-xq}
\[
X_q=X_{q-1}(v,-q),
\qquad
\dot u_q=p(v,-q)\quad(2.1.1).
\]

\VerdierSemanticAnchor{ch3d-2-3-2-verify}
Il résulte de (2.1.2) que $u_q$ est un quasi-isomorphisme, d'où la propriété
3). Il reste à vérifier la propriété 2), i.e. à vérifier que le noyau et
l'image du morphisme $L^{-q}\longrightarrow L^{-q+1}$ sont des éléments de
$M^\wedge$ ; ce qui résulte immédiatement des suites exactes :

\VerdierSemanticAnchor{ch3d-2-3-2-exact-sequences-2}
\[
0\longrightarrow\operatorname{Im}(L^{-q}\longrightarrow L^{-q+1})
\longrightarrow\operatorname{Ker}(L^{-q+1}\longrightarrow L^{-q+2})
\longrightarrow H^{-q+1}(X_q)\longrightarrow0,
\]
\[
0\longrightarrow\operatorname{Ker}(L^{-q}\longrightarrow L^{-q+1})
\longrightarrow L^{-q}\longrightarrow
\operatorname{Im}(L^{-q}\longrightarrow L^{-q+1})\longrightarrow0.
\]

\VerdierSemanticAnchor{ch3d-2-3-2-conclusion}
Ceci achève la preuve de la proposition.

\VerdierSemanticAnchor{ch3d-2-3-3}
\medskip\noindent\textbf{Proposition 2.3.3.}
\textit{Soient $M\subset M'$ deux ensembles d'objets de $\mathcal A$ tels que
$M$ possède les propriétés $(E_4)$ et $(E_5)$, que $M'$ possède les
propriétés $(E_1)$ et $(E_4)$ et que tout objet de $M'$ soit quotient d'un
objet de $M$. Soit $Y\in\operatorname{Ob}\operatorname{comp}(\mathcal A)$ un
complexe tel que :}

\noindent\textit{1) Les composants $Y^i$, $i\in\mathbf Z$, de $Y$ soient
des éléments de $M'$.}
\endgroup
